\section*{5. The Griffiths group}

\subsubsection*{5.0.}
Let $S$ be proper and smooth over an algebraically closed field $k$.
Denote by

\medskip
\noindent\textbf{(5.0.1)} $Z^i(S)$ the group of algebraic cycles of
codimension $i$ on $S$;

\noindent\textbf{(5.0.2)} $Z^i_{\mathrm{alg}}(S)$ the subgroup of
$Z^i(S)$ consisting of cycles algebraically equivalent to zero; and, for
every prime number $\ell\ne\operatorname{char}(k)$,

\noindent\textbf{(5.0.3)} $Z^i_{\ell\text{-}\mathrm{coh}}(S)$ the subgroup
of $Z^i(S)$ consisting of cycles whose cohomology class in
$H^{2i}(S,\mathbb Q_\ell(i))$ is zero.

It is conjectured in \textup{[4]} that this group is independent of $\ell$;
this is known at least when $\operatorname{char}(k)=0$.  Set
\begin{equation}\tag{5.0.4}
\mathrm{Grif}^i_\ell(S)
 =Z^i_{\ell\text{-}\mathrm{coh}}(S)/Z^i_{\mathrm{alg}}(S).
\end{equation}
If, moreover, $S$ is projective and is equipped with a projective embedding,
set

\medskip
\noindent\textbf{(5.0.5)} $Z^{i,\mathrm{prim}}(S)$ equal to the subgroup
of $Z^i(S)$ consisting of primitive cycles (that is, cycles whose
cohomology class is primitive), and

\noindent\textbf{(5.0.6)}
$\mathrm{Prim}^{2i}_{\mathrm{alg}}(S,\mathbb Q_\ell(i))$ equal to the
$\mathbb Q$-vector subspace of $H^{2i}(S,\mathbb Q_\ell(i))$ spanned by
the cohomology classes of primitive algebraic cycles.

\subsubsection*{Corollary 5.1.}
Under the hypotheses of Theorem~4.1, one has
\begin{equation}\tag{5.1.1}
\dim_{\mathbb Q}\!\left(
  \mathrm{Grif}^n_\ell(X_{\bar\eta})\otimes_{\mathbb Z}\mathbb Q
\right)
\geq
\dim_{\mathbb Q}\!\left(
  \mathrm{Prim}^{2n}_{\mathrm{alg}}(X,\mathbb Q_\ell(n))
\right).
\end{equation}
Indeed, by (4.1.5), the kernel of the composite map
\begin{equation*}
\mathbb Q\otimes_{\mathbb Z}Z^{n,\mathrm{prim}}(X)
 \xrightarrow{\ \mathrm{res.}\ \text{to }X_{\bar\eta}\ }
\mathbb Q\otimes_{\mathbb Z}Z^n_{\ell\text{-}\mathrm{coh}}
  (X_{\bar\eta})
 \longrightarrow \mathrm{Grif}^n_\ell(X_{\bar\eta})
\end{equation*}
is contained in
$Z^n_{\ell\text{-}\mathrm{coh}}(X)\otimes_{\mathbb Z}\mathbb Q$.
Consequently,
$\mathrm{Prim}^{2n}_{\mathrm{alg}}(X,\mathbb Q_\ell(n))$ is a quotient
of a subspace of
$\mathrm{Grif}^n_\ell(X_{\bar\eta})\otimes_{\mathbb Z}\mathbb Q$,
which proves the assertion.

\subsubsection*{Remark 5.2.}
For $X$ a quadric hypersurface in $\mathbb P^{2m+1}$, one has
\begin{equation*}
\dim_{\mathbb Q_\ell}\!\left(
  \mathrm{Prim}^{2m}_{\mathrm{alg}}(X,\mathbb Q_\ell(m))
\right)=1.
\end{equation*}
The print has twist $n$ here and omits the outer closing parenthesis; the
middle degree $2m$ forces the twist $m$.

In particular, when $m=2$, for $X$ a quadric in $\mathbb P^5$, the section
of $X$ by a general hypersurface of degree $>5$ contains a curve which is
$\mathbb Q_\ell$-homologous to zero but of which no nonzero multiple is
algebraically equivalent to zero.\footnote{The print inserts the literal
parenthetical ``(4.2)'' at this point.  Direct inspection of the native scan
at 9000-dpi-equivalent enlargement confirms the glyphs; the numerical
specialization from $\mathbb P^{2m+1}$ to $\mathbb P^5$ is $m=2$.}

\begin{thebibliography}{9}
\bibitem{GriffithsPeriodsXX}
P. Griffiths, ``On the Periods of Certain Rational Integrals III,''
forthcoming in \emph{Annals of Mathematics}.

\bibitem{GrothendieckBrauerIIIXX}
A. Grothendieck, ``Le Groupe de Brauer III'' (especially \S\S 9--10), in
\emph{Dix Exposés sur la Cohomologie des Schémas}, North-Holland, 1968.

\bibitem{JouanolouCohomologieXX}
J.-P. Jouanolou, \emph{Cohomologie $\ell$-adique}, thesis, Paris, 1969.

\bibitem{KleimanCyclesXX}
S. Kleiman, ``Algebraic Cycles and the Weil Conjectures,'' in
\emph{Dix Exposés sur la Cohomologie des Schémas}, North-Holland, 1968.

\bibitem{ManinCurvesXX}
J. Manin, ``Rational Points of Algebraic Curves over Function Fields,''
\emph{Translations of the American Mathematical Society}.
\end{thebibliography}
